\clearpage
\section*{References and source roles}
\addcontentsline{toc}{section}{References and source roles}
\begin{thebibliography}{99}

\bibitem{BohmSontacchi1978}
Corrado B\"ohm and Giovanna Sontacchi,
\emph{On the existence of cycles of given length in integer sequences like
$x_{n+1}=x_n/2$ if $x_n$ is even, and $x_{n+1}=3x_n+1$ otherwise},
\emph{Atti della Accademia Nazionale dei Lincei. Rendiconti}
\textbf{64} (1978), no.~3, 260--264.
The fixed-return formula is used at its printed scope; a return length may
have a proper divisor as its exact period.

\bibitem{Crandall1978}
Richard E. Crandall,
\emph{On the ``$3x+1$'' problem},
\emph{Mathematics of Computation} \textbf{32} (1978), no.~144,
1281--1292.
Section~7 is the published antecedent for the accelerated affine cycle
formula and minimum-orbit inequality.  The companion does not present its
small finite exclusion as a stronger historical cycle bound.

\bibitem{Lagarias1990}
Jeffrey C. Lagarias,
\emph{The set of rational cycles for the $3x+1$ problem},
\emph{Acta Arithmetica} \textbf{56} (1990), no.~1, 33--53,
doi:10.4064/aa-56-1-33-53.
Theorem~2.1, equations~(2.1), (2.4)--(2.7), and the primitive-cycle
distinction supply the controlling parity-word and rational-cycle lineage.

\bibitem{Stanley2024}
Richard P. Stanley,
\emph{Enumerative Combinatorics}, volume~2, second edition,
Cambridge Studies in Advanced Mathematics, volume~208,
Cambridge University Press, 2024,
doi:10.1017/9781009262538.
Exercise~7.89(a), equations~(7.149)--(7.150) on printed p.~480, and its
solution on printed p.~555 identify the multivariate Lyndon series with
primitive necklaces enumerated by the occurrences of each colour.

\bibitem{Urata2003}
Toshio Urata,
\emph{The Collatz problem over $2$-adic integers},
\emph{The Bulletin of Aichi University of Education}, Natural Science,
\textbf{52} (2003), 5--11.
Printed pp.~8--11 give the exponent numerator, periodic exponent sequences,
periodic rational points, and the explicit integrality question.

\bibitem{LaarhovenDeWeger2012}
Thijs Laarhoven and Benne de Weger,
\emph{The Collatz conjecture and De Bruijn graphs},
arXiv:1209.3495v1, 2012.
Source lines 303--321 provide the necklace/Lyndon-word and primitive
$3n+b$ comparison.  The companion retains the exact reduced-denominator
conditions supplied by Lagarias.

\bibitem{Tao2026v7}
Terence Tao,
\emph{Almost all orbits of the Collatz map attain almost bounded values},
arXiv:1909.03562v7, revision of 16 July 2026; published in
\emph{Forum of Mathematics, Pi} \textbf{10} (2022), e12,
doi:10.1017/fmp.2022.8.
The version-pinned source is used only for the exact finite affine offset and
the path-characterization statement cited in the text.

\bibitem{DrosteKuichVogler2009}
Manfred Droste, Werner Kuich, and Heiko Vogler, editors,
\emph{Handbook of Weighted Automata}, Springer, 2009,
ISBN 978-3-642-01491-8, eISBN 978-3-642-01492-5.
This is a framework source for weighted paths, semiring-valued series, and
finite weighted automata; no Collatz-specific result is attributed to it.

\bibitem{KuichSalomaa1986}
Werner Kuich and Arto Salomaa,
\emph{Semirings, Automata, Languages}, EATCS Monographs on Theoretical
Computer Science, volume 5, Springer, 1986.
This is a framework source for free monoids, semirings, formal series,
finite-support polynomials, and Cauchy multiplication.

\bibitem{Steinberg2016}
Benjamin Steinberg,
\emph{Representation Theory of Finite Monoids}, Universitext, Springer,
2016, ISBN 978-3-319-43930-3, eBook ISBN 978-3-319-43932-7.
The minimum-ideal result in the text is proved directly and then compared
with Proposition 13.1 and Lemma 13.2 of this source.

\bibitem{ConnesConsani2015}
Alain Connes and Caterina Consani,
\emph{Geometry of the arithmetic site},
arXiv:1502.05580, 2015; published in
\emph{Advances in Mathematics} \textbf{291} (2016), 274--329,
doi:10.1016/j.aim.2015.11.045.
The source is used only for its arithmetic-site square,
$\operatorname{Conv}(\mathbb Z\times\mathbb Z)$, quotient, and Frobenius
statements.  No Collatz-specific result is attributed to it.

\bibitem{BhattLurie2022}
Bhargav Bhatt and Jacob Lurie,
\emph{The prismatization of $p$-adic formal schemes},
arXiv:2201.06124, 2022.
The source is used for the definition and functoriality of the Cartier--Witt
stack and for the distinction between that functoriality and Witt-vector
Frobenius.

\bibitem{GrosLeStumQuiros2023v3}
Michel Gros, Bernard Le Stum, and Adolfo Quir\'os,
\emph{Absolute calculus and prismatic crystals on cyclotomic rings},
arXiv:2310.13790v3, 30 June 2025.
The source is used only for its stated one-variable cyclotomic equivalences
and hypotheses; it is not cited as a Collatz construction.

\bibitem{Sarig1999}
Omri M. Sarig,
\emph{Thermodynamic formalism for countable Markov shifts},
\emph{Ergodic Theory and Dynamical Systems} \textbf{19} (1999), no.~6,
1565--1593, doi:10.1017/S0143385799146820.
The source is used for the fixed-base periodic definition of Gurevich
pressure and its countable-shift hypothesis boundary.  The full-shift
partition calculation in the text is proved directly.

\bibitem{ConwayRadinSadun1999}
John H. Conway, Charles Radin, and Lorenzo Sadun,
\emph{On angles whose squared trigonometric functions are rational},
\emph{Discrete \& Computational Geometry} \textbf{22} (1999), no.~3,
321--332, doi:10.1007/PL00009463; arXiv:math-ph/9812019v1.
Theorems~1--2 and the $d=1$ St\o rmer analysis in Section~2 supply the
published Gaussian-prime coordinates and rational independence for fixed
rational-tangent angle relations.  They do not state the uniform Collatz
prefix-rank conjecture.

\bibitem{Calcut2009}
Jack S. Calcut,
\emph{Gaussian integers and arctangent identities for $\pi$},
\emph{The American Mathematical Monthly} \textbf{116} (2009), no.~6,
515--530, doi:10.1080/00029890.2009.11920967.
The Main Lemma and Corollaries~1--3 supply the exact $\mu_4$ restriction and
the published single-rational-angle non-torsion criterion used as a
crosswalk.  The companion gives an independent Gaussian-valuation proof for
its exceptional two-letter lattice.

\bibitem{GasullLucaVarona2023}
Armengol Gasull, Florian Luca, and Juan L. Varona,
\emph{Three essays on Machin's type formulas},
\emph{Indagationes Mathematicae} \textbf{34} (2023), no.~6, 1373--1396,
doi:10.1016/j.indag.2023.07.002; arXiv:2302.00154v2.
Theorem~1 classifies its stated two-term, nonzero-right-hand-side formulas
with $2$-integer arguments and includes the exceptional
$\arctan(1/2)+\tfrac12\arctan(3/4)=\pi/4$ identity.  Its scope does not
include the all-length Collatz prefix relation lattice.

\bibitem{Deaconu1995}
Valentin Deaconu,
\emph{Groupoids Associated with Endomorphisms},
\emph{Transactions of the American Mathematical Society} \textbf{347}
(1995), no.~5, 1779--1786.
The triple-groupoid formula is used as a literature prototype.  The
companion does not apply the paper's compact self-covering theorem to the
positive odd integers; its discrete topological assertions are proved
directly.

\bibitem{HuberWustholz2022}
Annette Huber and Gisbert W\"ustholz,
\emph{Transcendence and Linear Relations of $1$-Periods},
Cambridge Tracts in Mathematics, volume 227, Cambridge University Press,
2022.
Definition~8.1 supplies the exact commutative-square convention for
$1$-motive morphisms; Sections~10.2 and 12 supply the Kummer logarithm and
the comparison of $1$-periods with periods of $1$-motives.  The
Collatz-specific tori and characters are constructed in this companion.

\bibitem{HuberMullerStach2025}
Annette Huber and Stefan M\"uller-Stach, with contributions by Benjamin
Friedrich and Jonas von Wangenheim,
\emph{Periods and Nori Motives}, Springer, annotated version dated
1 July 2025.
Definition~9.1.1 supplies the two primitive edge types in the effective
diagram of pairs and fixes the contravariant direction of functoriality.
The Cayley and multiplication pair maps used here are proved explicitly.

\bibitem{BarbieriVialeKahn2016}
Luca Barbieri-Viale and Bruno Kahn,
\emph{On the Derived Category of $1$-Motives},
\emph{Ast\'erisque} \textbf{381} (2016); arXiv:1009.1900v2.
The source supplies the general definition, morphism, embedding, Albanese,
and Cartier-dual framework.  It does not state the point $(2,3)$ or any of
the path-indexed character squares proved in this companion.

\bibitem{Wirth2022}
Laura Wirth,
\emph{Weighted Automata, Formal Power Series and Weighted Logic},
BestMasters, Springer Spektrum, 2022,
doi:10.1007/978-3-658-39323-6.

\bibitem{BelagaMignotte2000}
Edward G. Belaga and Maurice Mignotte,
\emph{Cyclic Structure of Dynamical Systems Associated with $3x+d$
Extensions of Collatz Problem}, IRMA preprint 2000/018, 2000.
\url{https://www.researchgate.net/publication/2459588}.

\end{thebibliography}
